%%================================================
%% Chapter 4 - Evaluation Results
%%================================================
\chapter{ผลการทดลองและการประเมินผล}
\label{chapter4}

\section{ภาพรวมการทดลอง}
\label{sec:ch4_overview}
การประเมินผลระบบต้นแบบ WheelSense ในบทนี้จัดทำเพื่อดูว่าองค์ประกอบหลักที่พัฒนาไว้สามารถทำงานร่วมกันได้ในระดับต้นแบบหรือไม่ โดยแบ่งการทดสอบออกเป็นส่วนที่สอดคล้องกับสถาปัตยกรรมในบทที่~\ref{chapter3} ได้แก่ อุปกรณ์ sensor และอุปกรณ์ภาคสนาม การระบุตำแหน่งภายในอาคารด้วย RSSI/KNN การทำงานของ mobile gateway และ server การใช้งาน dashboard/HMI พร้อมการประเมิน SUS และการทดสอบ AI/MCP pipeline ในด้านความสอดคล้องของการเลือกเครื่องมือ ความสม่ำเสมอของคำตอบ และเวลาแฝงของแต่ละ layer

การทดลองทั้งหมดเป็นการทดสอบในพื้นที่ทดลองและสภาพแวดล้อมของโครงงาน ไม่ใช่การนำระบบไปใช้งานจริงในสถานดูแลผู้สูงอายุแบบต่อเนื่อง ดังนั้นผลที่รายงานจึงเป็นผลเพื่อยืนยันความเป็นไปได้ของระบบต้นแบบและเพื่อระบุข้อจำกัดที่ต้องพัฒนาต่อ ไม่ใช่การรับรองความพร้อมใช้งานทางการแพทย์หรือการใช้งานจริงในทุกสภาพแวดล้อม ตารางที่~\ref{tab:ch4_experiment_scope} แสดงขอบเขตการทดลองที่ใช้ในบทนี้

\begin{table}[!htbp]
    \centering
    \caption{ภาพรวมขอบเขตการทดลองของระบบ WheelSense}
    \label{tab:ch4_experiment_scope}
    \renewcommand{\arraystretch}{1.18}
    \small
    \setlength{\tabcolsep}{3pt}
    \begin{tabular}{|L{0.24\textwidth}|L{0.31\textwidth}|L{0.32\textwidth}|}
        \hline
        \textbf{ส่วนที่ทดสอบ} & \textbf{ประเด็นที่ประเมิน} & \textbf{หลักฐานที่ใช้ประกอบผล} \\ \hline
        Sensor และอุปกรณ์ภาคสนาม & ความถี่ IMU, ความครบถ้วนของ telemetry, การจำแนก motion state, การเชื่อมต่อ Polar & ตารางสรุป sensor, ตาราง motion state, ภาพอุปกรณ์ และผล mobile gateway \\ \hline
        Indoor localization & ผลของ RSSI fingerprint, KNN และ max-RSSI fallback เมื่อสัญญาณผันผวน & ตารางผล RSSI/KNN ในภาคผนวก และการวิเคราะห์ข้อจำกัดของสัญญาณ \\ \hline
        Mobile gateway และ server & BLE connection, MQTT/REST publish, reconnect, database latency และ dashboard update & ตาราง gateway, ตาราง server/database latency และภาพ runtime \\ \hline
        Dashboard/HMI และ SUS & การทำงานตามบทบาทผู้ใช้ task completion และ usability score ที่มีหลักฐานประกอบ & ตาราง task result และ usability summary โดยระบุสถานะของ SUS score อย่างระมัดระวัง \\ \hline
        AI/MCP pipeline & IRR, cosine similarity และ latency ของ Layer 1, 2, 4 และ 5 & ชุดทดสอบ AI 36 กรณี และการวัด latency 30 รอบ \\ \hline
    \end{tabular}
\end{table}

จากตารางที่~\ref{tab:ch4_experiment_scope} การรายงานผลจะไม่แยกเป็นคำอธิบายเชิงทฤษฎีซ้ำกับบทที่ 2 หรือบทที่ 3 แต่จะเน้นว่าทดสอบในเงื่อนไขใด ใช้วิธีใด ผลที่ได้เป็นอย่างไร และผลนั้นสะท้อนข้อจำกัดของระบบส่วนใด การจัดรูปแบบนี้ยึดแนวทางเดียวกับรายงานรุ่นพี่ที่อธิบายเงื่อนไขการทดลองก่อน แล้วจึงอ้างตารางหรือรูปเพื่อสรุปผลที่วัดได้
\FloatBarrier

\section{การทดสอบ Sensor และอุปกรณ์ภาคสนาม}
\label{sec:ch4_sensor_field}
การทดสอบส่วนนี้ตรวจสอบการทำงานของ M5StickC Plus2 ที่ติดตั้งกับรถเข็นเพื่ออ่านค่า IMU และ battery ร่วมกับ Polar Verity Sense ที่อ่านข้อมูลสัญญาณชีพผ่าน mobile gateway การทดสอบทำในพื้นที่ทดลองโดยกำหนดลักษณะการเคลื่อนไหวของรถเข็นเป็น 4 กลุ่ม ได้แก่ หยุดนิ่ง เคลื่อนที่ปกติ เลี้ยวแรง และเหตุการณ์ลักษณะกระแทก เพื่อดูว่าข้อมูล sensor มีความต่อเนื่องเพียงพอและสามารถใช้จำแนกสถานะเบื้องต้นได้หรือไม่

\begin{figure}[!htbp]
    \centering
    \includegraphics[width=0.68\textwidth]{chapter4/m5stickcinwheelchiar.png}
    \caption{การติดตั้ง M5StickC Plus2 บนรถเข็นสำหรับทดสอบข้อมูล IMU และ telemetry}
    \label{fig:ch4-m5-wheelchair-test}
\end{figure}

ขั้นตอนการทดลองเริ่มจากการให้ M5StickC Plus2 อ่านค่า IMU ด้วยรอบเวลาที่กำหนด จากนั้นส่งข้อมูล telemetry ผ่าน mobile gateway ไปยัง broker และ backend โดยทดสอบ 10 รอบ รอบละ 60 วินาที เป้าหมายการอ่านข้อมูลอยู่ที่ 20 Hz คิดเป็นข้อมูลที่คาดหวังประมาณ 12,000 packet และรับได้ 11,808 packet ส่วนข้อมูลการเคลื่อนไหวถูกแบ่งเป็นหน้าต่างข้อมูล 170 window เพื่อเปรียบเทียบกับการสังเกตระหว่างทดลอง ผลสรุปเชิงระบบแสดงในตารางที่~\ref{tab:ch4_sensor_summary}

\begin{table}[!htbp]
    \centering
    \caption{ผลสรุปการทดสอบอุปกรณ์ภาคสนาม}
    \label{tab:ch4_sensor_summary}
    \renewcommand{\arraystretch}{1.18}
    \small
    \setlength{\tabcolsep}{3pt}
    \begin{tabular}{|L{0.33\textwidth}|L{0.16\textwidth}|L{0.16\textwidth}|L{0.24\textwidth}|}
        \hline
        \textbf{ตัวชี้วัด} & \textbf{ค่าเฉลี่ย} & \textbf{P95} & \textbf{หมายเหตุ} \\ \hline
        IMU sampling rate & 19.6 Hz & 20.3 Hz & ทดสอบ 10 รอบ รอบละ 60 วินาที ใกล้เคียงเป้าหมาย 20 Hz \\ \hline
        telemetry completeness & 98.4\% & 99.1\% & รับได้ 11,808 จาก 12,000 packet ที่คาดหวัง \\ \hline
        device/gateway to broker latency & 41 ms & 156 ms & ทดสอบภายใต้เครือข่ายปกติของพื้นที่ทดลอง \\ \hline
        Polar connection success & 94.4\% & -- & สำเร็จ 34 จาก 36 ครั้ง พบการ re-pairing บางรอบ \\ \hline
    \end{tabular}
\end{table}

เมื่อพิจารณาข้อมูล motion state การจำแนกสถานะหยุดนิ่งและการเคลื่อนที่ปกติให้ผลดีกว่าสถานะเลี้ยวแรงและเหตุการณ์ลักษณะกระแทก เนื่องจากสองสถานะหลังเกิดในช่วงเวลาสั้นและขึ้นกับแรงที่ผู้ทดสอบกระทำต่อรถเข็น ตารางที่~\ref{tab:ch4_motion_detection} แสดงจำนวนหน้าต่างข้อมูลและผลที่ตรวจจับได้

\begin{table}[!htbp]
    \centering
    \caption{ผลการตรวจจับสถานะการเคลื่อนไหวจากหน้าต่างข้อมูล IMU}
    \label{tab:ch4_motion_detection}
    \renewcommand{\arraystretch}{1.18}
    \small
    \setlength{\tabcolsep}{3pt}
    \begin{tabular}{|L{0.25\textwidth}|L{0.14\textwidth}|L{0.15\textwidth}|L{0.13\textwidth}|L{0.18\textwidth}|}
        \hline
        \textbf{สถานะ} & \textbf{จำนวน window} & \textbf{ตรวจจับได้} & \textbf{Accuracy} & \textbf{หมายเหตุ} \\ \hline
        stationary & 50 & 49 & 0.98 & baseline ของการสั่นต่ำ \\ \hline
        normal motion & 50 & 47 & 0.94 & รถเข็นเคลื่อนที่ปกติ \\ \hline
        sharp turn & 40 & 36 & 0.90 & gyro magnitude สูงขึ้นในช่วงสั้น \\ \hline
        impact-like & 30 & 27 & 0.90 & ตรวจเทียบกับการสังเกตระหว่างทดลอง \\ \hline
    \end{tabular}
\end{table}

\begin{figure}[!htbp]
    \centering
    \includegraphics[width=0.72\textwidth]{chapter4/ch4-motion-accuracy.png}
    \caption{กราฟสรุป accuracy ของการตรวจจับสถานะการเคลื่อนไหวจาก IMU}
    \label{fig:ch4-motion-accuracy}
\end{figure}

ผลการทดสอบแสดงว่า M5StickC Plus2 สามารถอ่านและส่งข้อมูล IMU ได้ใกล้เคียงค่าที่ออกแบบไว้ ขณะที่ telemetry completeness อยู่ในระดับที่เพียงพอสำหรับระบบต้นแบบ อย่างไรก็ตามการจำแนกเหตุการณ์ที่เกิดเร็ว เช่น การเลี้ยวแรงหรือการกระแทก ยังต้องอาศัยข้อมูลระยะยาวและการทวนสอบกับเหตุการณ์จริงมากขึ้นก่อนใช้เป็นตัวชี้วัดความเสี่ยง

\subsection{ผลการประเมินความแม่นยำในการวัดระยะทาง ความเร็ว และความเร่ง}
การประเมินความแม่นยำของข้อมูลการเคลื่อนไหวดำเนินการในพื้นที่ทดสอบภายในอาคาร โดยติดตั้ง M5StickC Plus2 บนล้อรถเข็นและกำหนดเส้นทางตรงระยะอ้างอิง 10.00 เมตร ผู้ทดสอบจำนวน 2 คนเข็นรถผ่านเส้นทางเดียวกันคนละ 10 รอบ รวม 20 รอบการทดสอบ จากนั้นเปรียบเทียบระยะทางและความเร็วที่ระบบคำนวณได้กับระยะอ้างอิงและเวลาที่จับด้วยวิธี manual timing calibration ผลการทดสอบสรุปแสดงในตารางที่~\ref{tab:ch4_sensor_distance_speed}

\begin{table}[!htbp]
    \centering
    \caption[ผลการทดสอบระยะทางและความเร็ว]{ผลการทดสอบความแม่นยำในการวัดระยะทางและความเร็วของเซ็นเซอร์รถเข็นในอาคาร (จำนวน 20 รอบการทดสอบ)}
    \label{tab:ch4_sensor_distance_speed}
    \renewcommand{\arraystretch}{1.2}
    \small
    \begin{tabular}{|c|c|c|c|c|c|c|}
        \hline
        \textbf{ผู้ทดสอบ} & \shortstack{\textbf{ระยะทาง}\\\textbf{จริง (m)}} & \shortstack{\textbf{ระยะทาง}\\\textbf{ที่วัดได้ (m)}} & \shortstack{\textbf{Error}\\\textbf{(\%)}} & \shortstack{\textbf{ความเร็ว}\\\textbf{เฉลี่ยจริง}\\\textbf{(m/s)}} & \shortstack{\textbf{ความเร็ว}\\\textbf{ที่วัดได้}\\\textbf{(m/s)}} & \shortstack{\textbf{Error}\\\textbf{(\%)}} \\ \hline
        คนที่ 1 & 10.00 & 9.85 $\pm$ 0.12 & 1.50 & 1.20 & 1.17 $\pm$ 0.05 & 2.50 \\ \hline
        คนที่ 2 & 10.00 & 9.78 $\pm$ 0.14 & 2.20 & 1.35 & 1.31 $\pm$ 0.07 & 2.96 \\ \hline
        \textbf{ค่าเฉลี่ย} & \textbf{10.00} & \textbf{9.82 $\pm$ 0.13} & \textbf{1.85} & \textbf{1.28} & \textbf{1.24 $\pm$ 0.06} & \textbf{2.73} \\ \hline
    \end{tabular}
\end{table}

จากตารางที่~\ref{tab:ch4_sensor_distance_speed} พบว่าการคำนวณระยะทางมีค่าคลาดเคลื่อนเฉลี่ยร้อยละ 1.85 และการวัดความเร็วมีค่าคลาดเคลื่อนเฉลี่ยร้อยละ 2.73 ภายใต้เงื่อนไขการทดสอบในอาคารบนเส้นทางตรงระยะสั้น ผลดังกล่าวเพียงพอสำหรับใช้เป็นตัวชี้วัดเชิงพฤติกรรมระดับต้นแบบ เช่น การเคลื่อนที่ต่อเนื่อง ระยะที่เคลื่อนในช่วงเวลาหนึ่ง และการเปลี่ยนแปลงความเร็วโดยประมาณ แต่ยังไม่ควรตีความเป็นค่ามาตรวิทยาที่ใช้แทนเครื่องมือวัดมาตรฐาน

ในส่วนของการวัดความเร่ง สัญญาณ acceleration และ angular velocity ถูกใช้เป็น feature สำหรับแยกจังหวะหยุดนิ่ง เคลื่อนที่ เลี้ยวแรง และเหตุการณ์กระแทก ผลเชิงความเร่งจึงรายงานร่วมกับผล motion-window ในตารางที่~\ref{tab:ch4_sensor_acceleration_feature} แทนการสร้างตัวเลขใหม่ที่ไม่มีข้อมูลดิบรองรับ โดยใช้จำนวน window และผลตรวจจับจากการทดสอบ motion state ชุดเดียวกับตารางที่~\ref{tab:ch4_motion_detection}

\begin{table}[!htbp]
    \centering
    \caption[ผลการตรวจ feature ความเร่งจาก IMU]{ผลการตรวจ feature ความเร่งและ gyro จาก IMU สำหรับจำแนกสถานะการเคลื่อนไหว}
    \label{tab:ch4_sensor_acceleration_feature}
    \small
    \renewcommand{\arraystretch}{1.18}
    \begin{tabular}{|L{2.5cm}|C{1.8cm}|L{3.4cm}|C{2.0cm}|L{3.0cm}|}
        \hline
        \textbf{สถานะทดสอบ} & \textbf{จำนวน window} & \textbf{feature ที่ใช้ตรวจ} & \textbf{ผลตรวจจับ} & \textbf{ข้อสังเกต} \\ \hline
        หยุดนิ่ง & 50 & acceleration magnitude ต่ำร่วมกับ deadband & 49/50 & สะท้อน baseline noise ต่ำหลังกรองสัญญาณ \\ \hline
        เคลื่อนที่ปกติ & 50 & แนวโน้ม velocity และ acceleration ต่อเนื่อง & 47/50 & เหมาะกับการตรวจการเคลื่อนที่ระยะสั้น \\ \hline
        เลี้ยวแรง & 40 & gyro magnitude และการเปลี่ยนทิศทางเร็ว & 36/40 & พลาดได้เมื่อเหตุการณ์สั้นหรือ RSSI/IMU fluctuation สูง \\ \hline
        กระแทกคล้ายชน & 30 & acceleration magnitude peak และ jerk โดยประมาณ & 27/30 & ใช้เป็น alert candidate ได้ แต่ยังต้องยืนยันด้วยข้อมูลบริบท \\ \hline
    \end{tabular}
\end{table}

ผลในตารางที่~\ref{tab:ch4_sensor_acceleration_feature} แสดงว่า EMA, gyro offset calibration และ deadband threshold ที่อธิบายในบทที่~\ref{chapter3} ช่วยลดสัญญาณรบกวนช่วงหยุดนิ่งและทำให้ feature ของการเคลื่อนที่มีความต่อเนื่องมากขึ้น อย่างไรก็ตามเหตุการณ์ที่เกิดเร็ว เช่น เลี้ยวแรงหรือกระแทก ยังมีความคลาดเคลื่อนสูงกว่าสถานะหยุดนิ่งและเคลื่อนที่ปกติ จึงควรใช้เป็นสัญญาณช่วยเตือนเบื้องต้นร่วมกับตำแหน่ง เวลา และข้อมูลผู้ดูแล ไม่ใช่ผลตัดสินเหตุการณ์อันตรายโดยลำพัง

\FloatBarrier

\section{การทดสอบ Indoor Localization ด้วย KNN}
\label{sec:ch4_localization}
การทดสอบการระบุตำแหน่งภายในอาคารใช้ข้อมูล RSSI จาก BLE node ในพื้นที่ทดลอง โดย gateway สแกนสัญญาณและส่งรายการ RSSI พร้อม timestamp เข้าสู่ backend เพื่อประเมินตำแหน่งด้วย fingerprint และ KNN ในกรณีที่ fingerprint ยังไม่พร้อมหรือค่าความมั่นใจต่ำ ระบบสามารถใช้ max-RSSI เป็น baseline หรือ fallback ได้ หัวข้อนี้จึงอธิบายสภาพพื้นที่ทดลอง วิธีจัดชุดข้อมูล ผลที่ได้ และข้อสังเกตจากความผันผวนของ RSSI ต่อเนื่องในรูปแบบเดียวกับการทดสอบส่วนอื่นของบทนี้

\begin{figure}[!htbp]
    \centering
    \includegraphics[width=0.34\textwidth]{chapter4/user-location.png}
    \caption{ตัวอย่างหน้าจอ mobile gateway ที่ใช้ตรวจสถานะตำแหน่งและข้อมูลจาก BLE node}
    \label{fig:ch4-user-location}
\end{figure}

พื้นที่ทดลองจากชุดข้อมูล RSSI แบ่งห้องหนึ่งออกเป็น 4 zone ขนาดประมาณ \(4 \times 4\) เมตรต่อ zone โดยใช้ฉากผ้าเป็นขอบเขตระหว่างพื้นที่ และติดตั้ง ESP32-S3 BLE anchor จำนวน 4 ตัวไว้บริเวณกึ่งกลางของแต่ละ zone ที่ระดับความสูงประมาณ 0.7 เมตร อุปกรณ์ tag คือ M5StickC Plus2 ที่โฆษณา BLE ด้วยอัตราประมาณ 5 Hz หรือคาบ 200 ms และอยู่ที่ระดับความสูงประมาณ 0.7 เมตรเช่นกัน ข้อมูลหนึ่งตัวอย่างประกอบด้วย timestamp และค่า RSSI จาก anchor ทั้ง 4 ตัวในรูปแบบ \((timestamp, r_1, r_2, r_3, r_4)\)

ชุดทดลองแบ่งเป็น 3 กรณี กรณีละ 48 episode และใช้การแบ่งข้อมูลแบบเดียวกันคือ 13 episode สำหรับ train และ 35 episode สำหรับ test โดย Case 1 เป็นการตรวจ body orientation แบบ facing-in/facing-out, Case 2 เป็น boundary transition ระหว่างในและนอก zone และ Case 3 เป็น multi-zone trajectory ที่ต้องทำนายลำดับการเคลื่อนที่ผ่าน 4 zone วิธี KNN ใช้ \(k=5\) และ Euclidean distance สำหรับ Case 1 และ Case 2 จะทำนายจาก RSSI 4 มิติราย sample แล้วรวมผลเป็นระดับ episode ด้วยค่าเฉลี่ยของผลทำนาย ส่วน Case 3 แบ่ง episode เป็น 4 segment ต่อเนื่องและใช้ feature mean/std ของ RSSI จาก anchor ทั้ง 4 ตัวในแต่ละ segment เพื่อทำนายลำดับ zone

\begin{figure}[!htbp]
    \centering
    \includegraphics[width=0.86\textwidth]{chapter4/paperieee-combined-rssi-plot.pdf}
    \caption{ตัวอย่างสัญญาณ RSSI ต่อเนื่องจากชุดทดลอง Case 1, Case 2 และ Case 3}
    \label{fig:ch4_paperieee_rssi_plot}
\end{figure}

ผลการทดสอบ KNN จากชุดข้อมูลที่มีหลักฐานประกอบแสดงในตารางที่~\ref{tab:ch4_knn_result} โดยรายละเอียด raw data และ summary อยู่ในภาคผนวก~\ref{app:knn_summary} และแนบเอกสารเรื่อง \textit{Comparative Performance Evaluation of BLE-Based Indoor Motion Tracking Using Machine Learning and Large Language Models} ไว้ในภาคผนวก~\ref{app:paperieee_extract}

\begin{table}[!htbp]
    \centering
    \caption{ผลการทดสอบ RSSI/KNN จากชุดข้อมูลที่มีหลักฐานประกอบ}
    \label{tab:ch4_knn_result}
    \renewcommand{\arraystretch}{1.18}
    \small
    \setlength{\tabcolsep}{3pt}
    \begin{tabular}{|L{0.12\textwidth}|L{0.32\textwidth}|L{0.12\textwidth}|L{0.12\textwidth}|L{0.12\textwidth}|L{0.12\textwidth}|}
        \hline
        \textbf{Case} & \textbf{เงื่อนไขทดลอง} & \textbf{Accuracy} & \textbf{Precision} & \textbf{Recall} & \textbf{F1} \\ \hline
        Case 1 & Body orientation, facing-in/facing-out & 0.8857 & 0.8421 & 0.9412 & 0.8889 \\ \hline
        Case 2 & Boundary transition, in/out zone & 0.6857 & 0.6875 & 0.6471 & 0.6667 \\ \hline
        Case 3 & Multi-zone trajectory, exact sequence & 0.8857 & 0.8857 & 0.8857 & 0.8857 \\ \hline
    \end{tabular}
\end{table}

\begin{table}[!htbp]
    \centering
    \caption{ผลเปรียบเทียบ KNN กับ max-RSSI baseline}
    \label{tab:ch4_knn_baseline}
    \renewcommand{\arraystretch}{1.18}
    \small
    \setlength{\tabcolsep}{3pt}
    \begin{tabular}{|L{0.13\textwidth}|L{0.35\textwidth}|L{0.15\textwidth}|L{0.17\textwidth}|L{0.12\textwidth}|}
        \hline
        \textbf{Case} & \textbf{เงื่อนไขทดลอง} & \textbf{KNN} & \textbf{Max-RSSI} & \textbf{ส่วนต่าง} \\ \hline
        Case 1 & Body orientation, facing-in/facing-out & 0.8857 & 0.8000 & +0.0857 \\ \hline
        Case 2 & Boundary transition, in/out zone & 0.6857 & 0.6286 & +0.0571 \\ \hline
        Case 3 & Multi-zone trajectory, exact sequence & 0.8857 & 0.7714 & +0.1143 \\ \hline
    \end{tabular}
\end{table}

จากตารางที่~\ref{tab:ch4_knn_result} Case 1 และ Case 3 ให้ค่าความถูกต้องเท่ากันที่ 0.8857 โดย Case 3 ยังมีค่า per-position accuracy สูงกว่าค่า exact sequence เพราะการทำนายลำดับหนึ่ง episode ต้องถูกครบทุกตำแหน่งพร้อมกันจึงนับว่าถูก ส่วน Case 2 ลดลงเหลือ 0.6857 เนื่องจากตำแหน่งบริเวณ boundary มี RSSI จากหลาย anchor ใกล้เคียงกันและได้รับผลจาก body shadowing หรือ multipath พร้อมกันหลายช่องสัญญาณ KNN จึงสับสนได้ง่ายเมื่อ distribution ของ RSSI ซ้อนทับกัน ผลนี้สะท้อนว่า BLE RSSI ใช้ได้ในระดับ zone สำหรับ prototype แต่ยังต้องเก็บ fingerprint ในตำแหน่งจริงให้มากขึ้น ปรับจำนวน anchor หรือใช้ fallback เช่น max-RSSI/temporal smoothing เมื่อค่า confidence ต่ำ
\FloatBarrier

\section{การทดสอบ Mobile Gateway และ Server}
\label{sec:ch4_gateway_server}
ส่วน mobile gateway ถูกทดสอบเพื่อดูการรวมข้อมูลจาก M5StickC Plus2, Polar Verity Sense และข้อมูลจากโทรศัพท์ก่อนส่งเข้าสู่ platform ผ่าน MQTT หรือ REST API การทดสอบเน้นเวลาการเชื่อมต่อ BLE, latency จาก smartphone ไปถึง broker, อัตราความสำเร็จของ MQTT publish, ความครบถ้วนของ vitals telemetry และการทำงานเมื่อเกิด reconnect หรือ retry queue

\begin{figure}[!htbp]
    \centering
    \includegraphics[width=0.36\textwidth]{chapter4/PolarVerity-mobileapp.png}
    \caption{หน้าจอ mobile gateway สำหรับตรวจสอบการเชื่อมต่อ Polar Verity Sense และสถานะส่งข้อมูล}
    \label{fig:ch4-polar-mobile}
\end{figure}

ผลการทดสอบในตารางที่~\ref{tab:ch4_gateway_server} แสดงว่าเวลาการเชื่อมต่อ BLE GATT อยู่ในระดับไม่กี่วินาทีจากการทดสอบ 36 ครั้ง และการส่งข้อมูลจาก smartphone ไปถึง broker มีค่า p50 เท่ากับ 78 ms ส่วน MQTT publish ทดสอบ 1,200 ครั้งและส่งสำเร็จ 1,183 ครั้ง คิดเป็น 98.6\% ซึ่งเพียงพอสำหรับการทดสอบระบบต้นแบบ แต่ยังพบผลกระทบจากการ re-pairing และการสูญหายของ PPG block บางช่วง

\begin{table}[!htbp]
    \centering
    \caption{ผลการทดสอบ mobile gateway และ latency ของ server path}
    \label{tab:ch4_gateway_server}
    \renewcommand{\arraystretch}{1.16}
    \small
    \setlength{\tabcolsep}{2.5pt}
    \begin{tabular}{|L{0.28\textwidth}|L{0.13\textwidth}|L{0.13\textwidth}|L{0.34\textwidth}|}
        \hline
        \textbf{ตัวชี้วัด} & \textbf{P50/ค่าเฉลี่ย} & \textbf{P95} & \textbf{หมายเหตุ} \\ \hline
        BLE GATT connection time & 3.2 s & 5.8 s & ทดสอบการจับคู่ 36 ครั้งกับ M5StickC Plus2 และ Polar \\ \hline
        smartphone to broker latency & 78 ms & 215 ms & วัดจาก gateway ไปถึง broker \\ \hline
        MQTT connection success & 98.6\% & -- & สำเร็จ 1,183 จาก 1,200 publish มี retry/reconnect บางช่วง \\ \hline
        vitals completeness & 96.2\% & 98.9\% & รับได้ 2,309 จาก 2,400 block จาก Polar ผ่าน mobile gateway \\ \hline
        heart rate resting & 72.4 bpm & -- & ค่าที่ใช้เป็น baseline ระหว่างทดสอบ \\ \hline
        heart rate active & 95.6 bpm & -- & ค่าขณะเคลื่อนไหว/ทำกิจกรรม \\ \hline
        PPG block loss & 2.1\% & -- & บางช่วงหายจากการเคลื่อนไหวหรือ BLE reconnect \\ \hline
    \end{tabular}
\end{table}

สำหรับฝั่ง server การทดสอบตรวจสอบเส้นทาง ingest, validate, store และ query ผ่าน FastAPI, MQTT broker และ PostgreSQL/JSONB บน Raspberry Pi 5 โดยดูทั้งค่า latency ราย endpoint และการอ่านข้อมูลที่ dashboard หรือ AI ต้องใช้ซ้ำ

\begin{figure}[!htbp]
    \centering
    \includegraphics[width=0.72\textwidth]{chapter4/pi5-docker.png}
    \caption{ตัวอย่างการตรวจสอบบริการบน Raspberry Pi 5 และ Docker runtime ระหว่างการทดสอบ server}
    \label{fig:ch4-pi-docker}
\end{figure}

\begin{figure}[!htbp]
    \centering
    \includegraphics[width=0.72\textwidth]{chapter4/rasppi-monitor.png}
    \caption{หน้าจอ monitor ของ Raspberry Pi 5 ระหว่างรัน service และ container ของระบบ}
    \label{fig:ch4-rasppi-monitor}
\end{figure}

\begin{figure}[!htbp]
    \centering
    \includegraphics[width=0.72\textwidth]{chapter4/plot-mqtt-latency-box.png}
    \caption{กราฟสรุปการกระจายของ latency ในเส้นทาง MQTT ระหว่างการทดสอบ}
    \label{fig:ch4-mqtt-latency-plot}
\end{figure}

\begin{figure}[!htbp]
    \centering
    \includegraphics[width=0.72\textwidth]{chapter4/ch4-server-latency-summary.png}
    \caption{กราฟสรุปค่า P50 และ P95 ของ server path จากการทดสอบ ingest, API, alert และ room prediction}
    \label{fig:ch4-server-latency-summary}
\end{figure}

\begin{table}[!htbp]
    \centering
    \caption{ผลการทดสอบ latency ของ server path}
    \label{tab:ch4_server_latency}
    \renewcommand{\arraystretch}{1.16}
    \small
    \setlength{\tabcolsep}{3pt}
    \begin{tabular}{|L{0.30\textwidth}|L{0.14\textwidth}|L{0.14\textwidth}|L{0.14\textwidth}|L{0.18\textwidth}|}
        \hline
        \textbf{Server path} & \textbf{P50} & \textbf{P95} & \textbf{P99} & \textbf{หมายเหตุ} \\ \hline
        MQTT telemetry end-to-end & 46 ms & 178 ms & 304 ms & รับ telemetry จาก gateway \\ \hline
        REST API read latency & 63 ms & 212 ms & 341 ms & อ่านข้อมูลให้ dashboard \\ \hline
        Alert propagation latency & 71 ms & 236 ms & 388 ms & สร้างและส่งต่อ alert \\ \hline
        Room prediction publication & 58 ms & 191 ms & 325 ms & เผยแพร่ผลตำแหน่ง \\ \hline
    \end{tabular}
\end{table}

\begin{table}[!htbp]
    \centering
    \caption{ผลการทดสอบ latency ของ query ฐานข้อมูล}
    \label{tab:ch4_db_latency}
    \renewcommand{\arraystretch}{1.16}
    \small
    \setlength{\tabcolsep}{3pt}
    \begin{tabular}{|L{0.30\textwidth}|L{0.17\textwidth}|L{0.17\textwidth}|L{0.25\textwidth}|}
        \hline
        \textbf{Query path} & \textbf{P50} & \textbf{P95} & \textbf{หมายเหตุ} \\ \hline
        recent telemetry JSONB & 3 ms & 43 ms & อ่านค่าสถานะล่าสุดของผู้ใช้ \\ \hline
        alert list by patient state & 14 ms & 50 ms & อ่านรายการ alert สำหรับ dashboard \\ \hline
        recent sensor event & 12 ms & 35 ms & อ่านเหตุการณ์ล่าสุดจากอุปกรณ์ \\ \hline
        audit trail read & 24 ms & 51 ms & อ่านประวัติการทำงาน \\ \hline
    \end{tabular}
\end{table}

ผลจาก mobile gateway และ server สะท้อนว่าสถาปัตยกรรมที่ให้โทรศัพท์เป็น gateway หลักช่วยลดปัญหาข้อมูลซ้ำจากหลายแหล่งและทำให้สามารถตรวจสถานะการเชื่อมต่อได้ชัดขึ้น ขณะที่ backend มี latency ระดับเพียงพอสำหรับ dashboard ต้นแบบ อย่างไรก็ตาม BLE และ PPG stream ยังไวต่อการเคลื่อนไหวของผู้ใช้และการ reconnect จึงต้องเพิ่ม offline buffer และ diagnostic ที่ละเอียดขึ้นก่อนใช้งานระยะยาว
\FloatBarrier

\section{การประเมิน Dashboard/HMI และ SUS Score}
\label{sec:ch4_dashboard_hmi}
การประเมิน dashboard/HMI ใช้ task scenario ตามบทบาทผู้ใช้ เช่น Admin, Head Nurse, Supervisor, Observer และ Patient โดยมีผู้เข้าร่วมทดสอบ 15 คน และดูว่าผู้ใช้สามารถเข้าถึงข้อมูลผู้พักอาศัย สถานะอุปกรณ์ alert และหน้าที่เกี่ยวข้องกับ mobile gateway ได้ครบหรือไม่ ผลที่รายงานเป็นข้อมูลระดับต้นแบบและใช้เพื่อประเมินทิศทาง usability ของระบบ

\begin{table}[!htbp]
    \centering
    \caption{ผลสรุป usability ของ dashboard/HMI จากหลักฐานปัจจุบัน}
    \label{tab:ch4_hmi_usability}
    \renewcommand{\arraystretch}{1.18}
    \small
    \setlength{\tabcolsep}{3pt}
    \begin{tabular}{|L{0.28\textwidth}|L{0.17\textwidth}|L{0.45\textwidth}|}
        \hline
        \textbf{ตัวชี้วัด} & \textbf{ผลเฉลี่ย} & \textbf{การตีความ} \\ \hline
        task completion & 0.93 & ผู้เข้าร่วม 15 คนทำภารกิจหลักสำเร็จในระดับสูง \\ \hline
        average satisfaction & 4.31/5 & ความพึงพอใจต่อภาพรวมอยู่ในระดับดีสำหรับต้นแบบ \\ \hline
        alert readability & 4.42/5 & ข้อความและการจัดวาง alert อ่านได้ค่อนข้างชัดเจน \\ \hline
        role navigation clarity & 4.18/5 & การแยกบทบาทช่วยลดความสับสน แต่ยังต้องเพิ่ม workflow บางส่วน \\ \hline
        SUS mean / SD & 76.5 / 8.7 & สะท้อน usability ระดับใช้งานได้ดีในระบบต้นแบบ \\ \hline
    \end{tabular}
\end{table}

เพื่อให้เห็นงานย่อยในแต่ละบทบาท ตารางที่~\ref{tab:ch4_hmi_tasks} สรุป task ที่ใช้ทดสอบ dashboard และ mobile gateway ผลแสดงว่าหน้าที่ดูผู้พักอาศัยและอ่าน alert ทำได้ดีกว่างานที่เกี่ยวกับการแก้ไขอุปกรณ์หรือการตรวจ queue การส่งข้อมูล ซึ่งเป็นส่วนที่ควรพัฒนาต่อใน workflow ของผู้ดูแล

\begin{table}[!htbp]
    \centering
    \caption{ผลการทำภารกิจบน dashboard และ mobile gateway}
    \label{tab:ch4_hmi_tasks}
    \renewcommand{\arraystretch}{1.16}
    \small
    \setlength{\tabcolsep}{3pt}
    \begin{tabular}{|L{0.36\textwidth}|L{0.17\textwidth}|L{0.12\textwidth}|L{0.25\textwidth}|}
        \hline
        \textbf{ภารกิจ} & \textbf{completion} & \textbf{เวลาเฉลี่ย} & \textbf{หมายเหตุ} \\ \hline
        เปิดหน้า resident dashboard & 0.93 & 18 s & พบข้อมูลผู้พักอาศัยได้เร็ว \\ \hline
        acknowledge alert & 0.93 & 24 s & ต้องปรับการแยก priority ให้ชัดขึ้น \\ \hline
        ตรวจสถานะอุปกรณ์ & 0.87 & 29 s & บางรายการต้องเปิดหลายหน้า \\ \hline
        เปิดบริบทพื้นที่หรืออาคาร & 0.80 & 34 s & งานที่ใช้บริบทหลายหน้าใช้เวลามากกว่า \\ \hline
        อ่านสถานะ \texttt{/WheelSensor} & 0.87 & 21 s & หน้า gateway ช่วยแยกปัญหา BLE/MQTT ได้ \\ \hline
    \end{tabular}
\end{table}

\begin{figure}[!htbp]
    \centering
    \includegraphics[width=0.72\textwidth]{chapter4/ch4-hmi-task-completion.png}
    \caption{กราฟสรุป task completion ของ dashboard/HMI จากผู้เข้าร่วมทดสอบ 15 คน}
    \label{fig:ch4-hmi-task-completion}
\end{figure}

ผล usability แสดงว่า dashboard ช่วยให้ผู้ดูแลมองเห็นข้อมูลผู้พักอาศัยและสถานะอุปกรณ์ได้รวมศูนย์มากขึ้น โดยคะแนน SUS เฉลี่ย 76.5 อยู่ในระดับที่ใช้สนับสนุนความเหมาะสมของแนวทาง HMI สำหรับระบบต้นแบบ อย่างไรก็ตามงานที่เกี่ยวกับบริบทพื้นที่และการตรวจอุปกรณ์ยังใช้เวลามากกว่างานอ่านสถานะทั่วไป จึงควรปรับ navigation และตัวกรองข้อมูลในรอบพัฒนาถัดไป

\subsection{การทดสอบและการแสดงผลลัพธ์หน้าจอ HMI แยกตามบทบาทผู้ใช้}
เพื่อให้เห็นความสอดคล้องของการออกแบบหน้าจอแดชบอร์ดกับสิทธิ์ผู้ใช้ หัวข้อนี้นำเสนอภาพหน้าจอและผลการทดสอบแยกตาม 5 บทบาทหลัก ได้แก่ Admin, Head Nurse, Supervisor, Observer และ Patient โดยใช้ผู้เข้าร่วมทดสอบ 15 คนจากชุดเดียวกับตารางที่~\ref{tab:ch4_hmi_usability} และให้ผู้ทดสอบทำภารกิจที่สัมพันธ์กับบทบาทของตนเอง ส่วน EaseAI ถูกแยกเป็น cross-role assistant เพราะผู้ใช้หลายบทบาทสามารถเรียกใช้ได้ภายใต้ scope ที่แตกต่างกัน

\subsubsection{บทบาทผู้ดูแลระบบ (Admin Role)}
หน้าจอ Admin ใช้ตรวจภาพรวมของ workspace และความครบถ้วนของข้อมูลระบบ เช่น task, facility, floorplan, device, personnel, alert, message, support, settings และ EaseAI (รูปที่~\ref{fig:ch4_hmi_admin_contact}) ภารกิจที่ใช้ประเมินในบทบาทนี้เน้นการค้นหาอุปกรณ์ ตรวจสถานะผู้พักอาศัย และเปิดบริบทพื้นที่หรืออาคาร ซึ่งสัมพันธ์กับผลในตารางที่~\ref{tab:ch4_hmi_tasks} ที่งานตรวจสถานะอุปกรณ์และงานบริบทพื้นที่ใช้เวลามากกว่างานอ่านสถานะทั่วไป สะท้อนว่า Admin เหมาะกับการตรวจระบบรวม แต่ยังควรปรับตัวกรองและ navigation ในงานที่ต้องเปิดหลายหน้า

\begin{figure}[!htbp]
    \centering
    \includegraphics[width=0.98\textwidth]{chapter3/dashboard-hmi/role-admin-contact.png}
    \caption[หน้าจอ HMI บทบาท Admin]{ภาพหน้าจอระบบในมุมมอง Admin สำหรับตรวจภาพรวม workspace, task, device, facility, alert และ EaseAI}
    \label{fig:ch4_hmi_admin_contact}
\end{figure}
\FloatBarrier

\subsubsection{บทบาทหัวหน้าพยาบาล (Head Nurse Role)}
การทดสอบหน้าจอของ Head Nurse (รูปที่~\ref{fig:ch4_hmi_head_nurse_contact}) เน้นการติดตามผู้พักอาศัยหลายคนในแผนกและการจัดการภารกิจการดูแล ผู้ทดสอบต้องเปิดรายการผู้พักอาศัย อ่าน alert และยืนยันรายการที่ต้องส่งต่อ ผลการทำภารกิจ acknowledge alert สำเร็จเฉลี่ย 0.93 และใช้เวลาเฉลี่ย 24 วินาทีตามตารางที่~\ref{tab:ch4_hmi_tasks} แสดงว่าหน้าจอนี้สนับสนุนงานอ่านสถานะและส่งต่องานได้ค่อนข้างดี แต่ยังควรเพิ่มการแบ่ง priority ของ alert ให้ชัดขึ้น

\begin{figure}[!htbp]
    \centering
    \includegraphics[width=0.98\textwidth]{chapter3/dashboard-hmi/role-head-nurse-contact.png}
    \caption[หน้าจอ HMI บทบาท Head Nurse]{ภาพหน้าจอระบบในมุมมอง Head Nurse ที่ใช้ในการประเมิน alert, patient profile, task และ floorplan}
    \label{fig:ch4_hmi_head_nurse_contact}
\end{figure}
\FloatBarrier

\subsubsection{บทบาทผู้ควบคุมงาน (Supervisor Role)}
หน้าจอ Supervisor (รูปที่~\ref{fig:ch4_hmi_supervisor_contact}) ใช้สำหรับติดตามภาพรวมงานปฏิบัติการ เหตุการณ์ฉุกเฉิน งานที่ต้องส่งต่อ และข้อมูลบุคลากรหรือผู้พักอาศัยในขอบเขตที่ได้รับมอบหมาย ผู้ทดสอบสามารถมองเห็นงานค้างและบริบทพื้นที่ได้ แต่ภารกิจเปิดบริบทพื้นที่หรืออาคารมี completion เฉลี่ย 0.80 และใช้เวลาเฉลี่ย 34 วินาที จึงเป็นส่วนที่ควรปรับการจัดกลุ่มข้อมูลและการลิงก์ระหว่าง floorplan กับ task ให้เร็วขึ้น

\begin{figure}[!htbp]
    \centering
    \includegraphics[width=0.98\textwidth]{chapter3/dashboard-hmi/role-supervisor-contact.png}
    \caption[หน้าจอ HMI บทบาท Supervisor]{ภาพหน้าจอระบบในมุมมอง Supervisor สำหรับติดตามงาน เหตุการณ์ฉุกเฉิน บุคลากร และบริบทพื้นที่}
    \label{fig:ch4_hmi_supervisor_contact}
\end{figure}
\FloatBarrier

\subsubsection{บทบาทผู้สังเกตการณ์ (Observer Role)}
หน้าจอ Observer (รูปที่~\ref{fig:ch4_hmi_observer_contact}) ถูกจำกัดให้เป็นการอ่านข้อมูลและติดตามสถานะเป็นหลัก ผู้ใช้สามารถดู overview, alert, floorplan, task และ profile ที่ได้รับอนุญาต แต่ไม่ควรมี action ที่เปลี่ยนค่าคอนฟิกหลักของระบบ ผลคะแนน role navigation clarity เฉลี่ย 4.18/5 ในตารางที่~\ref{tab:ch4_hmi_usability} สะท้อนว่าการแยกบทบาทช่วยให้ผู้ใช้เข้าใจขอบเขตงานได้ แต่ยังควรเพิ่มข้อความสถานะหรือ disabled action เพื่อสื่อให้ชัดว่ารายการใดเป็น read-only

\begin{figure}[!htbp]
    \centering
    \includegraphics[width=0.98\textwidth]{chapter3/dashboard-hmi/role-observer-contact.png}
    \caption[หน้าจอ HMI บทบาท Observer]{ภาพหน้าจอระบบในมุมมอง Observer สำหรับติดตามสถานะและรายงานข้อมูลแบบ read-only}
    \label{fig:ch4_hmi_observer_contact}
\end{figure}
\FloatBarrier

\subsubsection{บทบาทผู้ป่วยและผู้พักอาศัยรายบุคคล (Patient Role)}
หน้าจอ Patient (รูปที่~\ref{fig:ch4_hmi_patient_contact}) ถูกประเมินเพื่อลดความสลับซับซ้อนของผู้ใช้งานสูงอายุหรือผู้พักอาศัย โดยจำกัดข้อมูลไว้เฉพาะระดับส่วนบุคคล เช่น ตารางกิจกรรมประจำวัน room controls, service request และ care plan ส่วนตัว ผู้ทดสอบสามารถเปิดหน้า resident dashboard ได้สำเร็จเฉลี่ย 0.93 ภายในเวลาเฉลี่ย 18 วินาที และไม่พบการเปิดเผยรายชื่อผู้พักอาศัยคนอื่นในหน้าจอของบทบาทนี้ จึงสอดคล้องกับแนวทาง role scope ที่อธิบายในบทที่~\ref{chapter3}

\begin{figure}[!htbp]
    \centering
    \includegraphics[width=0.98\textwidth]{chapter3/dashboard-hmi/role-patient-contact.png}
    \caption[หน้าจอ HMI บทบาท Patient]{ภาพหน้าจอระบบในมุมมอง Patient ครอบคลุม personal dashboard, schedule, room controls และช่องทางขอความช่วยเหลือ}
    \label{fig:ch4_hmi_patient_contact}
\end{figure}
\FloatBarrier

\subsubsection{การปฏิสัมพันธ์กับระบบผู้ช่วย AI (EaseAI Chat Assistant)}
ระบบผู้ช่วย EaseAI (รูปที่~\ref{fig:ch4_hmi_easeai}) ถูกประเมินในฐานะ assistant ข้ามบทบาท โดยให้ผู้ดูแลถามข้อมูลผู้พักอาศัยและเตรียม action candidate ที่ต้องผ่าน MCP execution gate ผู้ทดสอบให้คะแนนความพึงพอใจเฉลี่ย 4.31/5 และผลในตารางที่~\ref{tab:ch4_ai_execution_gate} ใช้ตรวจว่า action ที่มีผลเปลี่ยนสถานะยังต้องอยู่ภายใต้ขั้นตอน Propose $\rightarrow$ Confirm $\rightarrow$ Execute ไม่ใช่การปล่อยให้ LLM สั่งงานระบบจริงโดยลำพัง

\begin{figure}[!htbp]
    \centering
    \includegraphics[width=0.9\textwidth]{chapter3/dashboard-hmi/admin-easeai-chat.png}
    \caption[หน้าจอ EaseAI บน Dashboard]{ภาพหน้าจอการสนทนากับ EaseAI สำหรับสืบค้นข้อมูลและเตรียมคำสั่งผ่านระบบ MCP}
    \label{fig:ch4_hmi_easeai}
\end{figure}
\FloatBarrier

\section{การทดสอบ AI: Inter-Rater Reliability}
\label{sec:ch4_ai_irr}
การทดสอบ Inter-Rater Reliability ใช้เพื่อประเมินว่า AI/MCP pipeline เลือกเครื่องมือหรือแนวทางปฏิบัติได้สอดคล้องกับมนุษย์หรือไม่ ชุดคำสั่งทดสอบมีจำนวน 36 ข้อ แบ่งตามระดับความกำกวมเป็น 4 กลุ่ม โดยมี human raters จำนวน 2 คนเป็นผู้กำหนด ground truth ของเจตนาและเครื่องมือที่เหมาะสม ค่าที่รายงานเป็นสัดส่วนความสอดคล้องกับคำตอบอ้างอิง ไม่ใช่ Cohen's kappa

ผลการทดสอบแสดงในตารางที่~\ref{tab:ch4_ai_irr} ระบบทำได้ถูกต้อง 31 จาก 36 ข้อ หรือ 86.11\% โดยข้อผิดพลาดทั้ง 5 ข้ออยู่ในกลุ่มที่มีความกำกวมสูง

\begin{table}[!htbp]
    \centering
    \caption{ผลการทดสอบ AI Inter-Rater Reliability ตามระดับความกำกวม}
    \label{tab:ch4_ai_irr}
    \renewcommand{\arraystretch}{1.18}
    \small
    \setlength{\tabcolsep}{3pt}
    \begin{tabular}{|L{0.28\textwidth}|L{0.16\textwidth}|L{0.16\textwidth}|L{0.30\textwidth}|}
        \hline
        \textbf{กลุ่มคำสั่ง} & \textbf{จำนวนคำสั่ง} & \textbf{ถูกต้อง} & \textbf{หมายเหตุ} \\ \hline
        No-ambiguous & 10 & 10 & เจตนาชัดเจน \\ \hline
        Low-ambiguous & 10 & 10 & มีคำกำกวมเล็กน้อย \\ \hline
        Medium-ambiguous & 10 & 10 & ยังระบุ intent ได้ \\ \hline
        High-ambiguous & 6 & 1 & ระบบเลือกถามกลับหรือหยุดมากกว่าการเดา \\ \hline
        \textbf{รวม} & \textbf{36} & \textbf{31} & \textbf{86.11\%} \\ \hline
    \end{tabular}
\end{table}

\begin{table}[!htbp]
    \centering
    \caption{ผลการตรวจสอบ execution gate ของ AI/MCP pipeline}
    \label{tab:ch4_ai_execution_gate}
    \renewcommand{\arraystretch}{1.18}
    \small
    \setlength{\tabcolsep}{3pt}
    \begin{tabular}{|L{0.36\textwidth}|L{0.16\textwidth}|L{0.38\textwidth}|}
        \hline
        \textbf{ผลลัพธ์ของ gate} & \textbf{จำนวนกรณี} & \textbf{ความหมาย} \\ \hline
        Read-only answer & 13 & ตอบหรือสรุปข้อมูลโดยไม่สร้างคำสั่งเปลี่ยนสถานะ \\ \hline
        Execution candidate with confirmation & 18 & สร้าง candidate ได้ แต่ต้องรอผู้ใช้ยืนยันก่อนดำเนินการ \\ \hline
        Clarification / safety stop & 5 & ถามกลับหรือหยุดเมื่อคำสั่งกำกวมสูง \\ \hline
        Policy/schema block after validation & 0 & ไม่พบ candidate ที่ผ่าน Layer 4 แล้วถูก block เพิ่มในชุดทดสอบนี้ \\ \hline
        \textbf{รวม} & \textbf{36} & \textbf{ครบตามชุดคำสั่งทดสอบ} \\ \hline
    \end{tabular}
\end{table}

การที่ข้อผิดพลาดกระจุกตัวอยู่ในกลุ่ม High-ambiguous ไม่ได้หมายความว่าระบบสั่งงานผิดทั้งหมด แต่สะท้อนกลไก safety over guessing ของ Layer 1 Intent Router ที่ยอมถามกลับเมื่อความมั่นใจไม่พอ แทนที่จะเดาเจตนาแล้วส่งต่อไปยัง tool โดยตรง แนวทางนี้เหมาะกับระบบที่อาจเกี่ยวข้องกับผู้พักอาศัยและอุปกรณ์จริง เพราะลดโอกาสเกิดการสั่งงานผิดบริบท
\FloatBarrier

\section{การทดสอบ AI: Text-Distance / Cosine Similarity}
\label{sec:ch4_ai_cosine}
การทดสอบ Text-Distance ใช้ cosine similarity เพื่อประเมินความสม่ำเสมอของคำตอบจาก Layer 4 เมื่อได้รับคำถามเดิมซ้ำหลายรอบ การทดสอบเริ่ม session ใหม่ทุกครั้งและไม่ใช้ context window เดิม เพื่อให้เห็นผลของ retrieval ranking และลำดับบริบทที่ Layer 2 ส่งให้ Layer 4 โดยตรง

\begin{table}[!htbp]
    \centering
    \caption{ผลการทดสอบ Text-Distance / Cosine Similarity ของคำตอบ AI}
    \label{tab:ch4_ai_cosine}
    \renewcommand{\arraystretch}{1.18}
    \small
    \setlength{\tabcolsep}{3pt}
    \begin{tabular}{|L{0.30\textwidth}|L{0.20\textwidth}|L{0.40\textwidth}|}
        \hline
        \textbf{กลุ่มคำถาม} & \textbf{Cosine similarity} & \textbf{การตีความผล} \\ \hline
        General & 0.85 & คำตอบมีทิศทางค่อนข้างสม่ำเสมอ \\ \hline
        Health/context & 0.76 & คำตอบผันผวนมากกว่าเนื่องจากต้องพึ่งพาบริบทจาก RAG \\ \hline
    \end{tabular}
\end{table}

จากตารางที่~\ref{tab:ch4_ai_cosine} กลุ่ม General มีค่าเฉลี่ย 0.85 สูงกว่ากลุ่ม Health/context ที่ได้ 0.76 เหตุผลสำคัญคือคำถามด้านสุขภาพหรือบริบทผู้พักอาศัยต้องใช้ข้อมูลจาก RAG มากกว่า เมื่อลำดับเอกสารหรือ ranking ใน Layer 2 เปลี่ยนเล็กน้อย รูปแบบคำตอบของ LLM ใน Layer 4 จึงเปลี่ยนตาม แม้สาระหลักยังอยู่ในเรื่องเดียวกัน ผลนี้ชี้ว่าการปรับ retrieval ranking และ context packing จะช่วยให้คำตอบเสถียรมากขึ้น
\FloatBarrier

\section{การทดสอบ AI: Latency Analysis}
\label{sec:ch4_ai_latency}
การวัด latency ของ AI/MCP pipeline ใช้เพื่อระบุคอขวดหลังจากปรับสถาปัตยกรรมจากแนวทางเดิมที่ใส่กฎ เครื่องมือ และตัวอย่างจำนวนมากไว้ใน prompt เดียว ไปสู่ pipeline แบบแยกส่วน 5 layer การวัดครั้งนี้ตัดตัวแปรด้าน UI และเครือข่ายภายนอกออก แล้ววัดเฉพาะเส้นทาง blocking หลัก ได้แก่ Layer 1 Intent Router, Layer 2 Context Assembler, Layer 4 Constrained LLM และ Layer 5 Execution Gate จำนวน 30 รอบ

\begin{table}[!htbp]
    \centering
    \caption{ผลการทดสอบ latency ของ AI/MCP pipeline (n=30)}
    \label{tab:ch4_ai_latency}
    \renewcommand{\arraystretch}{1.18}
    \small
    \setlength{\tabcolsep}{3pt}
    \begin{tabular}{|L{0.34\textwidth}|L{0.20\textwidth}|L{0.20\textwidth}|L{0.18\textwidth}|}
        \hline
        \textbf{ส่วนที่วัด} & \textbf{P50} & \textbf{P95} & \textbf{หมายเหตุ} \\ \hline
        Layer 1 Intent Router & 198.40 ms & 248.05 ms & deterministic + semantic routing \\ \hline
        Layer 2 Context Assembler & 96.26 ms & 126.79 ms & retrieval และ context packing \\ \hline
        Layer 4 Constrained LLM & 14,257.39 ms & 36,487.60 ms & ส่วน inference หลัก \\ \hline
        Layer 5 Execution Gate & 40.38 ms & 78.62 ms & schema และ policy check \\ \hline
        \textbf{Total blocking path} & \textbf{14,592.43 ms} & \textbf{36,941.06 ms} & \textbf{L1 + L2 + L4 + L5} \\ \hline
    \end{tabular}
\end{table}

\begin{figure}[!htbp]
    \centering
    \includegraphics[width=0.76\textwidth]{chapter4/plot-llm-ttft-tokens.png}
    \caption{กราฟประกอบการวิเคราะห์เวลาแฝงของการประมวลผล LLM ใน pipeline}
    \label{fig:ch4-llm-latency-plot}
\end{figure}

ผลในตารางที่~\ref{tab:ch4_ai_latency} ชี้ว่าคอขวดหลักอยู่ที่ Layer 4 Constrained LLM เพราะเป็นขั้นตอนเดียวที่ต้องใช้การสร้างคำตอบหรือ candidate ผ่านโมเดลภาษา ขณะที่ Layer 1, Layer 2 และ Layer 5 ใช้เวลารวมกันระดับไม่กี่ร้อยมิลลิวินาที เมื่อเทียบกับบันทึกการพัฒนาแนวทางเดิมที่บางงานใช้เวลาประมาณ 30--45 นาที การแยก logic ออกจาก LLM ช่วยลด p50 ของเส้นทางหลักลงเหลือประมาณ 14.59 วินาที อย่างไรก็ตาม p95 ยังอยู่ที่ 36.94 วินาที จึงยังควรพัฒนา dynamic model multiplexing เพื่อไม่ให้คำสั่งง่ายต้องผ่านโมเดลขนาดเดียวกับงานที่ซับซ้อน
\FloatBarrier
